\section{Background}
\label{sec:background}

\noindent\textbf{LLM RL as token-level policy optimization.}
We follow the standard setup for RL fine-tuning of autoregressive language
models~\citep{ouyang2022training,grpo,dppo}. Given a prompt $x$, the policy $\tpi$
generates a response $y=(y_1,\dots,y_T)$ token by token.
We write $s_t=(x,y_1,\dots,y_{t-1})$ for the state at step $t$, so that
$\tpi(y\mid x)=\prod_{t=1}^{T}\tpi(y_t\mid s_t)$. A scalar reward $R(x,y)$ is
returned after the full response, and the objective is
$\mathcal{J}(\tpi)=\E_{y\sim\tpi}[R(x,y)]$. Responses are sampled from a \emph{behavior} (or rollout) policy $\bmu$,
typically a slightly stale copy of the current policy served by an optimized
inference engine.
The per-token importance ratio is
\begin{equation}
\label{eq:ratio}
r_t \;=\; \frac{\tpi(y_t\mid s_t)}{\bmu(y_t\mid s_t)}.
\end{equation}
Critic-free estimators such as GRPO~\citep{grpo} and its variants
\citep{ahmadian2024back,drgrpo} form the advantage $\Adv_t$ from group-relative
returns.
We treat $\Adv_t$ as given and inherit it unchanged from the underlying
algorithm.

\noindent\textbf{PPO and the ratio-based trust region.}
PPO~\citep{ppo} maximizes a clipped surrogate,
\begin{equation}
\label{eq:ppo}
L^{\mathrm{PPO}}(\tpi)=\E_{y\sim\bmu}\!\left[\sum_{t=1}^{|y|}
\min\!\Big(r_t\,\Adv_t,\;\operatorname{clip}(r_t,1-\eps,1+\eps)\,\Adv_t\Big)\right],
\end{equation}
which can equivalently be written as a masked policy-gradient objective
\begin{equation}
\label{eq:ppo-mask}
\begin{aligned}
&L^{\mathrm{PPO}}(\tpi)=\E_{y\sim\bmu}\!\left[\sum_{t=1}^{|y|} M_t^{\mathrm{PPO}}\,r_t\,\Adv_t\right],
\\[4pt]
&\text{where}\quad M_t^{\mathrm{PPO}}=
\begin{cases}
0, & \sign\big(\Adv_t\,(r_t-1)\big)>0 \;\text{ and }\; |r_t-1|>\eps,\\[2pt]
1, & \text{otherwise.}
\end{cases}
\end{aligned}
\end{equation}
The two factors of the mask $M_t^{\mathrm{PPO}}$ are the
\emph{proximity} criterion $|r_t-1|>\eps$ and the \emph{direction} criterion
$\sign\big(\Adv_t\,(r_t-1)\big)>0$, where $\sign(x)\in\{-1,0,+1\}$ returns the sign of $x$.
The proximity criterion measures how far the
training policy is from the behavior policy, but only through a single sampled token. Notably,
$D_{\mathrm{TV}}\big(\bmu(\cdot\mid s_t)\,\|\,\tpi(\cdot\mid s_t)\big)=
\tfrac12\,\E_{y_t\sim\bmu}\big[\,|r_t-1|\,\big]$,
which indicates that $|r_t-1|/2$ is a one-sample Monte Carlo estimate of the total variation distance.
Thus, PPO in effect gates on a noisy single-sample estimate of a
divergence~\citep{dppo}. 
The direction criterion encodes the asymmetric nature of trust-region control.
The mask acts only on outward-pointing updates.
Only updates that drive the policy \emph{further} from the behavior policy
threaten the trust region and need to be restrained.
Updates that pull the policy back toward it relax the trust-region constraint
and should be kept by the mask.
The sign of $\Adv_t(r_t-1)$ identifies whether the update is outward-pointing at
the sampled token.
It is positive exactly when $\Adv_t>0$ raises a probability already above the
behavior policy ($r_t>1$), or when $\Adv_t<0$ lowers a probability already below
it ($r_t<1$).
The PPO mask therefore removes only token updates that are both outside the trust
region and outward-pointing, while keeping corrective updates.

\noindent\textbf{DPPO and the divergence-based trust region.}
DPPO~\citep{dppo} keeps this masked structure and its asymmetric direction
criterion, but rebuilds the proximity criterion around the full divergence between
the two policies. Instead of gating on the single-sample estimate $|r_t-1|>\eps$, it
measures the per-token divergence directly,
$\Dtok_t=\Dtok\big(\bmu(\cdot\mid s_t)\,\|\,\tpi(\cdot\mid s_t)\big)$,
and uses $\Dtok_t>\delta$ as the proximity criterion.
This gives
\begin{equation}
\label{eq:dppo}
\begin{aligned}
&L^{\mathrm{DPPO}}(\tpi)=\E_{y\sim\bmu}\!\left[\sum_{t=1}^{|y|} M_t\,r_t\,\Adv_t\right],
\\[4pt]
&\text{where}\quad M_t=
\begin{cases}
0, & \sign\big(\Adv_t\,(r_t-1)\big)>0 \;\text{ and }\; \Dtok_t>\delta,\\[2pt]
1, & \text{otherwise.}
\end{cases}
\end{aligned}
\end{equation}
A common choice, and the one we use throughout, sets
$\Dtok=\KL(\bmu\,\|\,\tpi)$, the forward KL, which controls total variation
through Pinsker's inequality.

The full-vocabulary KL is not computable from a production rollout, which exposes
only the $K$ largest log-probabilities per step. The top-$K$ divergence addresses
this by restricting attention to a reduced support
$\Keep=\operatorname{TopK}(\bmu(\cdot\mid s_t),K)\cup\{y_t\}$, the $K$ highest-probability
tokens under the behavior policy augmented with the sampled token $y_t$ when it falls
outside the top $K$, and collapsing everything outside
$\Keep$ into a single aggregated tail bucket carrying the residual mass
$\bmu_{\mathrm{tail}}=1-\sum_{i\in\Keep}\bmu_i$ (and likewise $\tpi_{\mathrm{tail}}$). With this reduction,
\begin{equation}
\label{eq:topk-kl}
\KLtopk(\bmu\,\|\,\tpi)
\;=\; \sum_{i\in\Keep}\bmu_i\log\frac{\bmu_i}{\tpi_i}
\;+\; \bmu_{\mathrm{tail}}\log\frac{\bmu_{\mathrm{tail}}}{\tpi_{\mathrm{tail}}},
\end{equation}
which is a lower bound on the true KL (by the log-sum inequality) and, since
top-$K$ typically captures more than $99\%$ of the probability mass, an extremely
tight one~\citep{dppo}. Instantiating the mask of Eq.~(\ref{eq:dppo}) with this top-$K$ divergence, $\Dtok=\KLtopk$, gives the DPPO-TopK-KL method that we take as the starting point and primary baseline for this work.

DPPO upgrades only the proximity criterion.
The direction criterion is inherited from PPO unchanged, still flagging a token as
eligible for masking exactly when
$\sign\big(\Adv_t\,(r_t-1)\big)>0$.
The next section examines whether this sampled-ratio direction criterion tracks
the sign of the divergence change, and derives a replacement from the divergence
itself.
